% 基础语文试卷 - 展示材料、古诗、选择标记题型和自动跨页作文格
% 编译：xelatex 03-chinese-basic.tex

\documentclass{exam-zh}

\begin{document}

\title{高中语文练习卷}
\subject{语文}
\maketitle

\section{现代文阅读}

\begin{material}
  [
    title = {春日小记},
    author = {佚名},
    source = {（节选）}
  ]
  清晨，阳光越过窗台，照亮书桌上的一页纸。远处的树叶在风中轻轻摇动，
  新的一天就这样安静地开始了。
\end{material}

\begin{question}
  下列词语中，加点字读音正确的一项是
  \begin{select}[show-mark = true]
    \sitem 安静（jìn）
    \sitem* 摇动（yáo）
    \sitem 清晨（qīng）
  \end{select}
\end{question}

\section{古诗文阅读}

\begin{poem}[title = {登鹳雀楼}, author = {王之涣}]
  白日依山尽，\\
  黄河入海流。\\
  欲穷千里目，\\
  更上一层楼。
\end{poem}

\section{写作}

请以“窗外”为题，写一篇不少于 800 字的文章。

\compositiongrid
  [row-sep-ratio = 3, count-marks = {600, 800}]
  {1000}

\end{document}
